\section{Ablation studies}

% Pre-training setup:
% \begin{itemize}
% \item MoCo vs SimCLR (questions: is momentum useful? performance wrt queue size?)
% \item BERT vs no BERT
% \item Data augmentations (cropping, ICT, word masking/deleting, shuffling)
% \item CLS token vs averaging, training data (ccnet, trec-covid, etc), size of vocab
% \end{itemize}

In this section, we investigate the influence of different design choices of the method used in the previous section.
For these experiments, we use a reduced setting, where the model is pre-trained for 200,000 gradient steps, with a batch size of 2,048 (over 32 GPUs) on Wikipedia.
For the fine-tuning on MS MARCO, the model is trained on 8 GPUs with a batch size of $512$ for 20,000 gradient steps with the AdamW optimizer, without hard negatives.
%\eg{should we only report results after fine-tuning on MS MARCO? It will make the discussion simpler, and this is probably still the main use case of the paper}

%In this subsection we investigation the influence of different components of the method used in the previous sections. For this we use a minimal setting, where the model is pretrained for 200k gradient steps on 32 gpus, for the finetuning phase on MSMarco the model is trained using 8 gpus during 20k gradients steps without hard negatives. The hyperparameters used here can be found in Table~\ref{tab:hyperparameters}.

\begin{figure}[t]
    \centering
    \includegraphics{queue_size.pdf}
    \caption{\textbf{Queue size.} We report nDCG@10 before and after fine-tuning on MS MARCO.
    \aj{same for simclr?}}
    \label{fig:queue_size}
\end{figure}

\subsection{MoCo vs. SimCLR}
First, we compare the two contrastive pre-training methods: MoCo and SimCLR.
As in SimCLR, the number of negative examples is equal to the batch size, we train models with a batch size of \textcolor{red}{4,096} and restrict the queue in MoCo to the same number of elements.
Thus, in this experiments, we investigate the effect of using of momentum encoder for the keys, instead of backpropagating the gradient and using stochastic gradient descent.
We report results, with and without fine-tuning on MS MARCO in Table~\ref{tab:simclr_vs_moco}.
We observe that the difference between the two methods is not important, especially after training on MS MARCO.
Thus, we propose to use MoCo as it will allow to easily scale the number of negative examples used to train the model, an important ingredient for contrastive learning.
\eg{We need to re-run either SimCLR w/ batch size 2,048 or MoCo w/ batch size 4,096.}

\begin{table*}[h]
  \center
  \small
  \begin{tabular}{l ccccccccc}
    \toprule
           & NFCorpus  & NQ   & FiQA & ArguAna & Quora & DBPedia & SciDocs & FEVER & AVG \\
    \midrule
    %MoCo   & 25.9 & 13.1 & 12.9 & 32.8 & 69.0 & 18.9 & 12.2 & 56.8 & 29.9 \\
    %SimCLR & 24.2 & 21.6 & 13.0 & 33.7 & 74.9 & 17.9 & 13.6 & 56.1 & 31.9 \\
    %\midrule
    MoCo   & 31.1 & 37.2 & 23.4 & 41.6 & 82.7 & 36.9 & 15.7 & 71.5 & 41.5 \\
    SimCLR & 29.6 & 36.5 & 23.3 & 38.6 & 83.6 & 34.4 & 16.2 & 72.3 & 41.8 \\
    \bottomrule
  \end{tabular}
  \label{tab:simclr_vs_moco}
  \caption{\textbf{MoCo vs. SimCLR.} In this table, we report nDCG@10 on the BEIR benchmark for SimCLR and MoCo, after fine-tuning on the MS MARCO dataset.}
\end{table*}

\subsection{Queue size}
Next, we study the influence of the number of negatives used in the contrastive loss, by varying the queue size of the MoCo algorithm.
We consider values ranging from 2,048 to 131,072, and report results in Figure~\ref{fig:queue_size}.
We see that on average over the BEIR benchmark, increasing the number of negatives leads to better retrieval performance, especially in the unsupervised setting.
However, it is interesting to note that this effect is not equally strong for all datasets.

%\gi{Would be better with something similar with simclr?}

%The use of a queue without backpropagation allows us to easily scale the size of the pool of negatives. As shown in Table~\ref{tab:queue_size}, scaling the queue size improves both the zero-shot performance and the performance after finetuning on MSMarco.

\subsection{Data augmentations}
Third, we compare different ways to generate pairs of positive examples from a single document or chunk of text.
In particular, we compare random cropping, which leads to pairs with overlap, and the inverse cloze task, which was previously considered to pre-train retrievers.
Interestingly, as shown in Table~\ref{tab:data_augmentation}, the random cropping strategy strongly outperforms the inverse cloze task.
We argue that overlap between the two views of the data is beneficial, as it allows the network to learn lexical matching, similarly to BM25.
We also investigate whether additional data perturbations, such as random word deletion or replacement, are beneficial for retrieval.
After fine-tuning on MS MARCO, it seems that such data perturbations slightly degrade the performance of dense retrievers.

\begin{table*}[h]
  \center
  \small
  \begin{tabular}{l ccccccccc}
    \toprule
    & NFCorpus & NQ & ArguAna & Quora & DBPedia & SciDocs & FEVER & Overall \\
    \midrule
    ICT             & 30.5 & 37.5 & 32.6 & 82.3 & 33.1 & 14.6 & 67.9 & 39.5 \\
    Crop            & 31.5 & 44.1 & 39.5 & 84.2 & 39.6 & 16.3 & 74.8 & 44.3 \\ % 46005732 - moco - qsize 65536 - ratio 10-50% 256 tokens
    Crop + delete   & 31.6 & 42.3 & 39.2 & 84.6 & 39.7 & 16.2 & 71.8 & 43.4 \\ % 46211625 - moco - delete 10% - qsize 65536 - ratio 10-50% 256 tokens
    Crop + replace  & 30.5 & 43.2 & 42.6 & 84.6 & 39.2 & 16.0 & 73.8 & 43.6 \\ % 46211626 - moco - replace 10% - qsize 65536 - ratio 10-50% 256 tokens
    \bottomrule
  \end{tabular}
  \caption{\textbf{Data augmentions.} We report nDCG@10 after fine-tuning on MS MARCO.}
  \label{tab:data_augmentation}
\end{table*}


%Data augmentations are key component of contrastive learning methods used in computer vision, in the main results presented in this paper we do a very minimal use of them. In Table~\ref{tab:augmentations} we consider different data augmentation strategies. Namely we present results with random deletion, random token replacement and random masking of p\% of the tokens. Additionally we consider cropping with and without overlap. 

\subsection{Training data}
Finally, we study the impact of the pre-training data on the performance of our retriever, by training on Wikipedia, CCNet or a mix of both sources of data.
We report results in Table~\ref{tab:training_data}, and observe that there is no clear winner between the two data sources.
Unsurprisingly, pre-training on Wikipedia leads to better performance on datasets for which documents also come from Wikipedia, such as NaturalQuestions or FEVER.
On the other hand, on datasets from different domains than Wikipedia such as FiQA or ArguAna, training on the more diverse CCNet data leads to better results.
To get the best of both worlds, we consider to strategies to mix the two data sources.
In the 50/50\% strategy, examples are sampled uniformly across domain, meaning that half the batches are from Wikipedia and the other half from CCNet.
In the uniform strategy, examples are sampled uniformly over the union of the dataset.
Since CCNet is significantly larger than Wikipedia, this means that most of the batches are from CCNet.

\begin{table*}[h]
  \center
  \small
  \begin{tabular}{l cccccccccccccc}
    \toprule
           & NFCorpus & NQ & FiQA & ArguAna & Quora & DBPedia & SciDocs & FEVER & Overall \\
    \midrule
    %Wiki          & 27.6 & 17.7 & 35.6 & 75.4 & 21.0 & 13.3 & 64.5 & 63.0 & 33.0 \\ % 46005732 - moco - qsize 65536 - ratio 10-50% 256 tokens
    %CCNet         & 29.5 & 25.8 & 35.2 & 80.6 & 20.5 & 14.9 & 60.9 & 59.4 & 34.9 \\ % 46307600 - moco - cc-net - qsize 65536 - m=0.999
    %50/50\%       & 27.1 & 17.8 & 39.1 & 76.5 & 19.9 & 13.9 & 66.3 & 66.0 & 34.7 \\ % 46403749 - moco - 50-50 cc-net/wiki - lr 7e-5 - qsize 65536 - m=0.999
    %Uniform       & 29.1 & 19.2 & 36.2 & 79.0 & 22.6 & 13.6 & 64.0 & 62.7 & 35.6 \\ % 46435971 - moco - uniform cc-net/wiki - lr 7e-5 - qsize 65536 - m=0.999
    %\midrule
    Wiki          & 31.5 & 44.1 & 26.8 & 39.5 & 84.2 & 39.6 & 16.3 & 74.8 & 44.3 \\
    CCNet         & 31.8 & 39.9 & 31.6 & 43.3 & 84.5 & 36.1 & 16.0 & 66.4 & 42.4 \\
    50/50\%       & 31.7 & 43.2 & 27.1 & 35.9 & 84.7 & 39.5 & 15.8 & 77.4 & 44.0 \\
    Uniform       & 32.0 & 41.5 & 30.5 & 39.9 & 84.9 & 39.6 & 16.2 & 72.0 & 44.1 \\
    \bottomrule
  \end{tabular}
  \caption{\textbf{Training data.} We report nDCG@10 after fine-tuning on MS MARCO.}
  \label{tab:training_data}
\end{table*}

